\documentclass[11pt]{article}

% =====================================================================
% Identity Was All You Needed: Wallet-as-Identity, Sealed Disclosure,
% and Capability-Grant Access for Agents. Fourth in the Unbrowse
% paper series; academic format mirrors the trilogy. Every [shipped]
% claim is anchored in paper/anchors.tsv; the honest-gaps section is
% load-bearing, not decoration.
% =====================================================================

\usepackage[margin=1in]{geometry}
\usepackage{hyperref}
\usepackage{amsmath}
\usepackage{amssymb}
\usepackage{enumitem}
\usepackage{booktabs}
\usepackage{xcolor}
\hypersetup{colorlinks=true, linkcolor=blue, urlcolor=blue, citecolor=blue}

\title{\textbf{You Are Your Keys (Sorry)}\\
\large Wallet-as-Identity, Sealed Disclosure, and Capability-Grant Access for Agents}

\author{Lewis Tham\\ Unbrowse AI \\ \texttt{lewis@unbrowse.ai}
  \and Cayden Chik\\ Unbrowse AI \\ \texttt{cayden@unbrowse.ai}}
\date{June 2026}

\begin{document}
\maketitle

\begin{abstract}
An agent that signs every layer it touches~\cite{crypto} has solved \emph{how} an
action is attested, but not the two questions that decide whether it should run at
all: \emph{who} is asking, and \emph{who may see the answer}. This paper makes the
narrow claim that those two questions collapse into one act --- a signature --- and
need no account system, no bearer-token vault, and no per-resource ACL service to
answer. Concretely: (i) a
wallet signature is a first-class identity --- the same Ed25519 key~\cite{rfc8032}
that signs requests \emph{is} the principal, with no web2 account required beneath
it; (ii) every cached value is sealed to the wallet and revealed only under
signature, fail-closed, so disclosure is a cryptographic consequence of identity
rather than a policy check; and (iii) read access between principals is a signed,
scoped, revocable capability grant in the object-capability tradition~\cite{ocap,
macaroons}, where the granter's signature is the sole authority and there is no
ambient permission. Authentication, selective disclosure, and authorisation are
three readings of one signed object. We report what ships as running code, and we
are explicit about the standard guarantees this design does \emph{not} provide
(revocation latency, key loss, metadata exposure), because an identity paper that
claims to have closed those is not to be trusted.
\end{abstract}

\section{Introduction: who, and who may see}

The companion security paper establishes that one Ed25519 key can sign every layer
an agent descends through --- screen, browser, CLI, OS, kernel, packet --- so that
no action crosses a boundary unsigned~\cite{crypto}. That is a statement about
\emph{integrity of action}. It leaves open two orthogonal questions that every
multi-agent system must answer before an action is even admissible:

\begin{enumerate}[leftmargin=1.6em,itemsep=1pt]
  \item \textbf{Authentication} --- who is this principal? Conventionally answered by
  an account record and a credential the principal presents.
  \item \textbf{Authorisation and disclosure} --- may this principal read this
  resource? Conventionally answered by an access-control list consulted at read time,
  and a secret store that hands out the bytes once the check passes.
\end{enumerate}

The conventional answers each introduce a trusted third surface: an account
database, an ACL service, a secret vault. Each is a place where authority is
\emph{asserted} rather than \emph{proven}, and therefore a place an adversary can
route around or a misconfiguration can leak. This paper's claim is that for an agent
that already carries a signing key, all three surfaces disappear into the key
itself. Identity is the public key; disclosure is decryption under the private key;
authorisation is a signature by the resource owner naming who may decrypt. There is
nothing left to consult.

This is not a new cryptographic primitive --- it is a composition of well-understood
ones (digital signatures, authenticated encryption, object capabilities) arranged so
that the answer to ``who, and who may see'' is read off a signed object instead of
fetched from a service. The contribution is the arrangement and the demonstration
that it ships.

\section{Related work}

\paragraph{Wallet / key-as-identity.} Public-key-as-principal is the oldest idea in
the field: SPKI/SDSI binds authority to keys rather than names, and the
object-capability literature makes the key the unforgeable handle~\cite{ocap}. Our
contribution is not the idea but its operational shape for agents: the signing key
the agent uses for action integrity is reused, unchanged, as the authentication
principal, so there is no separate enrolment step and no account to phish.

\paragraph{Selective disclosure.} Anonymous-credential systems let a holder reveal
attributes without revealing more~\cite{camlys}; sealed-sender and end-to-end
encryption bind plaintext to a key. We use the simplest instance of the family:
authenticated encryption keyed by a value derived from the wallet, so a stored
record is opaque to everyone but the holder, and opening it re-verifies the
content-address. Disclosure is not a policy decision taken by a server; it is the
ability to decrypt.

\paragraph{Capabilities over ACLs.} The object-capability model replaces ``check the
subject against a list'' with ``hold an unforgeable, attenuable token''~\cite{ocap};
macaroons make such tokens delegatable and caveat-scoped with only a keyed
MAC~\cite{macaroons}. Our grants are the signed instance: a capability is a row
signed by the granter, naming a grantee (or a lineage pattern), a scope, and an
expiry; verification is signature-check plus scope-match, with no central authority
in the loop.

\paragraph{Transparency.} That every grant and every settled claim is appended to a
hash-chained, append-only log is the Certificate-Transparency
discipline~\cite{rfc6962} applied to authorisation: revocation and issuance are
themselves logged events, so the access history is auditable rather than implicit in
a mutable table.

\section{Method: one signed object, three readings}

\subsection{Authentication --- the signature is the principal}
A request carries an Ed25519 signature over its canonical body. Verification
recovers the public key; \emph{that key is the identity}. There is no lookup that
can fail open: an unsigned or mis-signed request has no principal and is rejected. A
key may optionally be bound to a web2 account for convenience, but the binding is
additive --- absent it, the principal is simply \texttt{wallet:$\langle$pubkey$\rangle$},
a fully usable standalone identity. \textbf{A wallet signature is a first-class
identity}; the account, where present, is a label on top of it, never the root of
it. \emph{[shipped]}

\subsection{Disclosure --- sealed to the key, opened by the key}
Every value the agent stores is sealed under authenticated encryption with a key
derived from the wallet secret, and the stored record holds only a commitment (a
content hash) plus opaque ciphertext. A value is \textbf{revealed only under
signature, fail-closed}: presenting the wrong key fails the authentication tag and
returns nothing; presenting the right key decrypts and then re-checks the
content-address before the plaintext is trusted. Disclosure is therefore a property
of \emph{holding the key}, not of passing a server-side permission check --- there
is no read endpoint that, misconfigured, hands the bytes to the wrong caller.
\emph{[shipped]}

\subsection{Authorisation --- a signed, scoped, revocable capability grant}
For one principal to read another's data, the resource owner signs a \textbf{scoped,
revocable capability grant}: a record naming the grantee (an exact key or a lineage
pattern), the scope it authorises, an expiry, and the granter's signature as the
sole authority. A read is admitted only if a grant exists whose signature verifies
against the granter's key, whose scope matches the request, and which has not
expired or been revoked; otherwise it fails closed. There is no ambient authority
and no central policy server --- the grant \emph{is} the permission, and revocation
is an appended event that the next verification observes. \emph{[shipped]}

\subsection{Why the three collapse into one}
Each reading is the same object viewed from a different angle. The principal is a
public key. Disclosure is decryption under its private half. Authorisation is a
signature \emph{by} a key naming another key. Authentication, selective disclosure,
and access control are not three subsystems that must agree; they are three uses of
the one signing identity, which is why none of them needs a trusted third surface to
adjudicate.

\section{Honest gaps --- what this does not solve}
An identity design that claims to have closed the hard problems is lying; we name
the open ones as plainly as the wins.

\begin{table}[h]
\centering
\begin{tabular}{lll}
\toprule
Concern & Status & Honest reading \\
\midrule
Revocation latency & open & a revoke is an appended event; a verifier that has not
yet observed it can still admit \\
Key loss / recovery & open & losing the private key loses the identity and every
sealed value under it \\
Metadata exposure & open & ciphertext hides values, not the existence/shape of the
access graph \\
Grant explosion & partial & lineage-pattern grants bound the count, but broad
patterns trade precision for convenience \\
\bottomrule
\end{tabular}
\caption{Open problems this design does not close. Each is a standard hard problem
in capability systems; naming them is the point.}
\label{tab:gaps}
\end{table}

These are not incidental. Revocation in a log-based capability system is
eventually-consistent by construction; key loss is the cost of having no account to
reset against; metadata privacy needs machinery (mixing, padding) this design does
not include. The contribution is the collapse of the three positive questions into
one signed object, \emph{not} a claim to have also solved revocation, recovery, and
traffic analysis.

\section{Conclusion}
Once an agent carries a signing key for the integrity of its actions, that key is
already enough to answer who is acting and who may see the result: identity is the
public key, disclosure is decryption under its private half, and authorisation is a
signed, scoped, revocable grant naming another key. Authentication, selective
disclosure, and access control collapse into three readings of one signed object,
removing the account database, the ACL service, and the secret vault as trusted
surfaces. We ship the three readings as running code and we decline to claim the
hard residue --- revocation latency, key loss, metadata exposure --- is solved. The
line between the collapse (real) and the residue (open) is the honest core of the
result.

\begin{thebibliography}{99}
\bibitem{crypto} L.~Tham, C.~Chik. \emph{Crypto Was All You Needed: A Signed,
Layer-Descending Stack for Agent Computation}. Unbrowse AI, 2026.
\bibitem{wedge} L.~Tham et al. \emph{Internal APIs Are All You Need}.
arXiv:2604.00694, 2026.
\bibitem{rfc8032} S.~Josefsson, I.~Liusvaara. \emph{Edwards-Curve Digital Signature
Algorithm (EdDSA)}. RFC 8032, IETF, 2017.
\bibitem{ocap} M.~S.~Miller. \emph{Robust Composition: Towards a Unified Approach to
Access Control and Concurrency Control}. PhD thesis, Johns Hopkins University, 2006.
\bibitem{macaroons} A.~Birgisson, J.~G.~Politz, U.~Erlingsson, A.~Taly, M.~Vrable,
M.~Lentczner. \emph{Macaroons: Cookies with Contextual Caveats for Decentralized
Authorization in the Cloud}. NDSS, 2014.
\bibitem{camlys} J.~Camenisch, A.~Lysyanskaya. \emph{An Efficient System for
Non-transferable Anonymous Credentials with Optional Anonymity Revocation}.
EUROCRYPT, 2001.
\bibitem{rfc6962} B.~Laurie, A.~Langley, E.~Kasper. \emph{Certificate Transparency}.
RFC 6962, IETF, 2013.
\end{thebibliography}

\end{document}
